\section{全球模型公司与 AI 产品公司格局}

前面讨论产品壁垒，本章进入公司格局。2025 年 Q1 的公司图谱里，OpenAI、Anthropic、Google、DeepSeek、ByteDance、Mira/Thinking Machines、SSI、Cursor、Manus 等角色被放在同一张竞争图中。广密的讨论带有强主观判断，但它可以帮助我们理解不同类型公司在 AI 价值链中的位置。

\begin{figure}[H]
\centering
\includegraphics[width=0.92\textwidth]{model-company-portfolio.png}
\caption{全球模型公司格局：不同公司押注模型、产品、生态、开源和资本结构。自制概念图，依据 01:15:11--01:54:32 对谈内容整理。}
\end{figure}

\begin{knowledgebox}{读图：不要把所有 AI 公司放在同一赛道}
Anthropic 更偏企业/Coding，OpenAI 同时押模型和产品，ByteDance 有产品和流量，DeepSeek 代表开源效率，Cursor/Manus 是应用执行层。不同位置的成功条件完全不同，不能只用同一个 benchmark 排名。
\end{knowledgebox}

\subsection{GPT-4.5、GPT-5 与 OpenAI 风险}

本节梳理 OpenAI 讨论。访谈提到 GPT-4.5 是否领先、GPT-5 为什么跳票、OpenAI 是否有失败风险、OpenAI 与微软关系、以及对 Anthropic MCP 协议的支持。课程化整理时应把这些看成几个变量：底座模型节奏、reasoning model 节奏、产品用户增长、云合作结构、协议生态和组织稳定性。

\begin{warningbox}{不要用单个发布判断一家公司}
OpenAI 的风险和优势都很大：它有最强品牌、用户入口和产品心智，也有组织复杂化、云关系、研究节奏和商业化压力。单个模型版本不能决定长期胜负。
\end{warningbox}

\subsection{DeepSeek、开源效率与中国信号}

本节看 DeepSeek。访谈把 DeepSeek 放在 Q1 的明星位置，强调它对全球模型格局和开源效率叙事的冲击。它的价值不只是“模型便宜”，而是说明中国团队可以通过工程效率、架构选择、训练 recipe 和开源传播影响全球认知。广密甚至用投资组合语言表达对 DeepSeek 的高权重偏好，这反映了他对开源效率路线的重视。

\begin{importantbox}{DeepSeek 的信号意义}
DeepSeek 让行业重新评估中国团队在 frontier 模型、训练效率、开源生态和全球开发者影响力上的位置。它不是单纯的价格故事，而是能力、成本、开放和传播叠加的故事。
\end{importantbox}

\subsection{Manus、Perplexity、Cursor：执行力强的产品公司}

本节看应用公司。访谈里 Manus、Perplexity、Cursor 被称为“模型盗火者”，甚至被调侃为“套壳之王”。这个说法表面上尖锐，但真正重点在执行力：它们把模型公司尚未产品化的能力快速转成用户体验。Perplexity 把搜索/研究重新组织，Cursor 把 Coding 变成高频工作流，Manus 把通用 Agent 的体验推到公众面前。

\begin{knowledgebox}{术语消化：模型公司、应用公司、基础设施公司}
\noindent\begin{tabularx}{\textwidth}{p{0.20\textwidth}p{0.34\textwidth}X}
\toprule
公司类型 & 核心资产 & 主要风险 \\
\midrule
模型公司 & base model、训练 infra、研究人才、API/平台 & 训练成本高、组织复杂、产品压力大。 \\
应用公司 & 用户入口、工作流、任务状态、反馈数据 & 被模型公司复制或被成本压垮。 \\
基础设施公司 & 推理、训练、环境、工具协议、云资源 & 需要跟随模型和应用需求变化。 \\
开源模型公司 & 成本效率、社区、透明度、生态扩散 & 商业化和持续训练资金压力。 \\
\bottomrule
\end{tabularx}
\end{knowledgebox}

\subsection{MCP、协议和生态位}

本节补足协议层。访谈提到 OpenAI 支持 Anthropic 的 MCP 协议，这说明 Agent 生态不只竞争模型，还竞争工具连接方式、上下文传递方式和安全边界。协议如果成为事实标准，会把模型、工具、企业系统和开发者生态连接起来。模型公司既要防止被协议商品化，又要避免自己被生态孤立。

\begin{importantbox}{Agent 生态的下一层竞争是协议}
浏览器、IDE、企业 SaaS、文件系统和数据库都可能变成 Agent 工具。谁定义模型如何安全地读取上下文、调用工具和返回结果，谁就在定义 Agent 生态的交通规则。
\end{importantbox}

\subsection{本章小结}

全球 AI 公司格局不能只按“谁的模型最强”排序。底座模型、开源效率、应用执行、协议生态、云合作和资本结构都在重排位置。2025 Q1 的真正变化，是模型能力开始更明确地流向 Coding、Agent 和工作流产品。

